% Part: incompleteness
% Chapter: introduction
% Section: historical-background

\documentclass[../../../include/open-logic-section]{subfiles}

\begin{document}

\olfileid[fr]{inc}{int}{bgr}

\olsection{Repères historiques}

Nous retracerons brièvement dans cette section quelques évolutions
historiques qui aideront à situer les théorèmes d'incomplétude. Nous
esquisserons d'abord l'histoire de la logique mathématique, puis nous
dirons quelques mots de celle des fondements des mathématiques.

\begin{digress}
  L'expression « logique mathématique » est ambiguë. L'adjectif
  « mathématique » peut qualifier l'objet étudié, comme dans « la
  logique des mathématiques », c'est-à-dire les principes du
  raisonnement mathématique; il peut aussi qualifier la méthode,
  comme dans « une étude mathématique de la logique », c'est-à-dire
  l'étude mathématique des principes du raisonnement. Le récit qui
  suit concerne la logique mathématique dans les deux sens, souvent
  à la fois.
\end{digress}

L'étude de la logique a commencé, pour l'essentiel, avec Aristote,
qui vécut approximativement de 384 à 322 avant notre ère. Ses
\emph{Catégories}, ses \emph{Premiers Analytiques} et ses
\emph{Seconds Analytiques} comprennent des études systématiques des
principes du raisonnement scientifique, dont une étude approfondie
et méthodique du syllogisme.

La logique d'Aristote a dominé la philosophie scolastique tout au
long du Moyen Âge. Au XVIIIe siècle encore, Kant soutenait qu'elle
était parfaite et n'avait pas besoin d'être révisée. Pourtant, la
théorie du syllogisme est beaucoup trop limitée pour rendre compte
d'autre chose que des aspects les plus superficiels du raisonnement
mathématique. Un siècle auparavant, Leibniz, contemporain de Newton,
avait imaginé un « calcul » complet du raisonnement logique et fait
quelques premiers pas vers sa conception, décrivant en substance une
forme de logique propositionnelle.

Le XIXe siècle a marqué un tournant pour la logique. En 1854, George
Boole publia \emph{The Laws of Thought}, une étude algébrique
approfondie de la logique propositionnelle qui n'est guère éloignée
des présentations modernes. En 1879, Gottlob Frege publia sa
\emph{Begriffsschrift} (« Écriture conceptuelle »), qui étend la
logique propositionnelle par les quantificateurs et les relations,
et comprend donc la logique du premier ordre. En réalité, les
systèmes logiques de Frege englobaient aussi la logique d'ordre
supérieur, entre autres. Dans ses \emph{Grundgesetze der Arithmetik}
(« Lois fondamentales de l'arithmétique »), Frege entreprit de
montrer que toute l'arithmétique pouvait être déduite, dans sa
\emph{Begriffsschrift}, d'hypothèses purement logiques. Ces
hypothèses se révélèrent malheureusement contradictoires, comme
Russell le signala à Frege en 1902. Mais, si l'on met de côté l'axiome en
cause, Frege a pour ainsi dire inventé à lui seul la logique
moderne, ce qui est un exploit saisissant. La logique des
quantificateurs fut aussi développée indépendamment, dans le
sillage de Boole, par des penseurs tournés vers l'algèbre, dont
Peirce et Schr\"oder.

Passons maintenant à l'histoire des fondements des mathématiques.
Puisque la logique joue un rôle important en mathématiques, ces
développements recoupent largement ceux que nous venons de décrire.
Frege, par exemple, élabora sa logique dans le but explicite de
montrer que toutes les mathématiques pouvaient reposer sur son seul
cadre logique. Il voulait notamment établir que les mathématiques
consistent en vérités \emph{analytiques} a priori, et non en
vérités \emph{synthétiques} a priori, comme le soutenait Kant.

On situe souvent la naissance des mathématiques proprement dites
chez les Grecs. Les \emph{Éléments} d'Euclide, rédigés vers 300 avant
notre ère, offrent déjà un exemple abouti des mathématiques
grecques, soucieuses de rigueur et de précision. Les définitions et
les preuves des \emph{Éléments} se retrouvent encore, à peu de
chose près, dans les manuels scolaires de géométrie (dans la mesure
où la géométrie est toujours enseignée à l'école). Ce modèle de
raisonnement mathématique a été tenu pour un paradigme de
l'argumentation rigoureuse, en mathématiques comme dans certaines
branches de la philosophie. Spinoza alla jusqu'à présenter des
arguments moraux et religieux à la manière d'Euclide, ce qui est
surprenant à lire!

Newton et Leibniz inventèrent le calcul infinitésimal au XVIIe
siècle. Une âpre querelle de priorité dura des siècles, mais la
plupart des spécialistes estiment aujourd'hui que leurs travaux
furent, pour l'essentiel, indépendants. Le calcul infinitésimal
fait intervenir, par exemple, des sommes infinies de quantités
infiniment petites. Ces caractéristiques nourrirent les critiques
de l'évêque Berkeley, selon qui croire en Dieu n'était pas moins
rationnel que de croire aux mathématiques de son temps. Les méthodes
du calcul infinitésimal furent largement employées au XVIIIe
siècle, notamment par Leonhard Euler, qui obtint des résultats
spectaculaires grâce à des calculs sur les sommes infinies.

Au XIXe siècle, les mathématiciens cherchèrent à répondre aux
critiques de Berkeley en donnant au calcul infinitésimal des
fondements plus solides. Les travaux de Cauchy, Weierstrass,
Bolzano et d'autres conduisirent aux définitions actuelles des
limites, de la continuité, de la dérivation et de l'intégration en
termes d'« epsilon et delta », sans aucune référence aux
infinitésimaux. Plus tard dans le siècle, les mathématiciens
tentèrent d'aller plus loin en expliquant tous les aspects du
calcul infinitésimal, y compris les nombres réels eux-mêmes, à
partir des nombres naturels. Comme le disait Kronecker : « Dieu a
créé les nombres entiers; tout le reste est l'œuvre de l'homme. »
En 1872, Dedekind rédigea « Continuité et nombres irrationnels »,
où il montrait comment « construire » les nombres réels comme des
ensembles de nombres rationnels (lesquels, comme vous le savez,
peuvent être vus comme des couples de nombres naturels). En 1888,
il écrivit \emph{Was sind und was sollen die Zahlen?} (« Que sont
les nombres naturels et que doivent-ils être? », approximativement),
dans le but d'expliquer les nombres naturels en termes purement
« logiques ». En 1887, Kronecker écrivit \emph{\"Uber den
Zahlbegriff} (« Sur le concept de nombre »), où il parlait de
représenter tous les objets mathématiques au moyen des entiers. En
1889, Giuseppe Peano donna des axiomes formels et symboliques pour
les nombres naturels.

La fin du XIXe siècle apporta aussi une audace nouvelle dans le
traitement de l'infini. Jusqu'alors, on abordait avec prudence les
objets et les structures infinis, comme l'ensemble des nombres
naturels : « une infinité de » signifiait « autant que l'on veut »
et « tendre vers une limite », « s'en approcher autant que l'on
veut ». Georg Cantor montra qu'on pouvait au contraire prendre
l'infini au pied de la lettre. Les travaux de Cantor, de Dedekind
et d'autres contribuèrent à installer la conception ensembliste
générale des mathématiques qui est aujourd'hui largement admise.

Nous arrivons ainsi aux développements de la logique et des
fondements au XXe siècle. En 1902, Russell signala à Frege le paradoxe
présent dans son système logique.\footnote{Précision éditoriale : l'année~1902 renvoie ici à la lettre du~16~juin adressée à Frege ; la source la donne comme date de découverte. Voir la transcription conservée par les Archives Bertrand Russell, BRACERS~53946, \url{https://bracers.mcmaster.ca/53946}.} En 1904, Zermelo démontra
le principe du bon ordre de Cantor en utilisant le « principe du
choix », dit aussi « axiome du choix »; la légitimité de cet axiome
suscita de nombreux débats. Entre 1910 et 1913 parurent les trois
volumes des \emph{Principia Mathematica} de Russell et Whitehead,
qui poursuivaient le programme de Frege visant à fonder les
mathématiques sur la logique. Mais Russell et Whitehead durent
adopter deux principes qu'il semblait difficile de justifier comme
purement logiques : un axiome de l'infini et un axiome de
« réductibilité ». Dans les années 1900, Poincar\'e critiqua
l'emploi des « définitions imprédicatives » en mathématiques. Dans
les années 1910, Brouwer commença à proposer de refonder toutes
les mathématiques sur une base « intuitionniste » qui évitait le
recours au principe du tiers exclu ($!A \lor \lnot !A$).

Quelle époque étrange!{} Le programme de réduction de toutes les
mathématiques à la logique porte aujourd'hui le nom de
« logicisme » et est généralement considéré comme ayant échoué,
en raison des difficultés déjà mentionnées. Celui qui consiste à
développer les mathématiques à partir de constructions mentales
intuitionnistes s'appelle « intuitionnisme »; on estime qu'il
impose des restrictions excessives à la pratique ordinaire des
mathématiques. Au tournant du siècle, David Hilbert, l'un des
mathématiciens les plus influents de tous les temps, défendait
ardemment les nouvelles méthodes abstraites introduites par Cantor
et Dedekind : « Personne ne nous chassera du paradis que Cantor a
créé pour nous. » Il était néanmoins sensible aux critiques des
fondements adressées à ces méthodes, aujourd'hui appelées, assez
curieusement, « classiques ». Il proposa de concilier les deux
exigences :
\begin{enumerate}
\item Représenter les méthodes classiques par des axiomes et des
  règles formelles; représenter les questions mathématiques par des
  !!{formula}s dans un système axiomatique.
\item Employer des méthodes sûres, « finitistes », pour démontrer
  la cohérence de ces systèmes déductifs formels.
\end{enumerate}

Les travaux de Hilbert contribuèrent largement à réaliser le premier
objectif. Dès 1899, il l'avait fait pour la géométrie dans son
ouvrage célèbre, \emph{Fondements de la géométrie}. Dans les années
suivantes, plusieurs de ses étudiants et collaborateurs travaillèrent
avec lui à étendre à d'autres domaines des mathématiques ce qu'il
avait fait pour la géométrie. Hilbert lui-même proposa des systèmes
d'axiomes pour l'arithmétique et l'analyse. Zermelo axiomatisa la
théorie des ensembles; Fraenkel, Skolem, von Neumann et d'autres
développèrent ensuite son travail. Vers le milieu des années 1920,
deux approches revendiquaient le titre d'axiomatisation de « toutes »
les mathématiques : les \emph{Principia Mathematica} de Russell et
Whitehead, et ce qu'on en vint à appeler la théorie des ensembles
de Zermelo--Fraenkel.

En 1921, Hilbert lança un programme de recherche visant à démontrer
la cohérence de ces systèmes. Il fut aidé par plusieurs de ses
étudiants, notamment Bernays, Ackermann et, plus tard, Gentzen.
L'idée fondamentale était de reformuler la possibilité d'une
!!{derivation} menant à une contradiction mathématique comme un
problème combinatoire portant sur des suites de symboles. Il
s'agissait de déterminer s'il pouvait exister une suite d'énoncés
constituant une !!{derivation} correcte, à partir des axiomes de
l'arithmétique, de l'analyse ou de la théorie des ensembles, d'une
contradiction telle que $!A \land \lnot !A$. Une preuve qu'une telle
suite de symboles est impossible serait elle-même une preuve
mathématique et pourrait donc être formalisée dans ces systèmes
axiomatiques. Autrement dit, il existerait un énoncé $\OCon$
affirmant, par exemple, que l'arithmétique est cohérente. De plus,
cet énoncé devrait être démontrable dans les systèmes en question,
surtout si sa preuve ne requiert que des moyens « finitistes » très
restreints.

L'objectif de faire en sorte que les systèmes
axiomatiques élaborés tranchent toute question mathématique peut
être précisé de deux manières. D'abord : pour tout énoncé~$!A$ du
langage d'un système axiomatique pour les mathématiques, ou bien
$!A$, ou bien $\lnot !A$ est démontrable à partir des axiomes.
Alors, chaque énoncé pourrait être démontré ou réfuté à partir
des axiomes, et aucune question ne resterait sans réponse.
Un système axiomatique qui possède cette propriété est dit
\emph{complet}. Bien sûr, pour un énoncé donné, déterminer
laquelle des deux possibilités se réalise pourrait rester difficile.
Mais il devrait exister, en principe, une méthode pour le faire.
En fait, pour les systèmes d'axiomes et de !!{derivation}s considérés
par Hilbert, la complétude entraînerait l'existence d'une telle
méthode, bien que Hilbert ne l'ait pas compris. La seconde manière
de comprendre cet objectif consiste à poser une exigence plus forte :
qu'il existe une méthode mécanique, calculatoire, permettant de
déterminer si un énoncé~$!A$ donné est ou non déductible des
axiomes.

En 1931, G\"odel démontra les deux « théorèmes d'incomplétude »,
qui montrèrent que ce programme ne pouvait aboutir. Sous les
hypothèses des théorèmes, aucun système axiomatique des mathématiques
ne peut être complet. En particulier, l'énoncé exprimant la
cohérence des axiomes ne peut être démontré dans le système
cohérent concerné; sa non-réfutabilité demande des hypothèses
supplémentaires.\footnote{Précision éditoriale : ces résultats concernent les théories effectives ayant la force arithmétique requise. Le second théorème porte sur leur énoncé usuel de cohérence et exige des conditions de démontrabilité supplémentaires ; il n'affirme pas, sous la seule cohérence, que cet énoncé est aussi irréfutable.}

Ce résultat porta un coup fatal au programme initial de Hilbert.
Mais, comme il arrive souvent en mathématiques, il ouvrit aussi
des voies de recherche passionnantes. S'il n'existe pas un unique
système formel embrassant toutes les mathématiques, il est naturel
d'élaborer des systèmes plus circonscrits et d'étudier ce qu'on
peut y démontrer. Il est également naturel de développer des
méthodes de preuve moins restrictives pour établir leur cohérence,
et de chercher à mesurer la difficulté de cette preuve. Puisque
G\"odel montra que les systèmes formels satisfaisant les hypothèses
d'incomplétude ont des questions qu'ils ne peuvent trancher, il est pertinent de chercher
des questions « intéressantes » qu'un système donné ne tranche pas
et de déterminer la force requise d'un système pour les résoudre.
Les logiciens poursuivent encore aujourd'hui ces recherches au
sein d'une nouvelle discipline mathématique : la théorie de la
démonstration.

\end{document}
